%-----------------------------------------------------------
%   Capítulo 3 - Análise de Requisitos
%-----------------------------------------------------------
\chapter{Análise de Requisitos}
\label{cap:requisitos}

Os requisitos seguintes resultam dos objetivos específicos da Introdução,
dos cenários CA01--CA03 (Secção~\ref{sec:cenarios-aplicacao}), das
conclusões do Estado da Arte (Secção~\ref{sec:sintese-estado-arte}) e do
levantamento e validação de requisitos junto dos docentes responsáveis
pelo projeto. Este processo confirma o
acesso bidirecional do sistema central, através de Ethernet, a CAN,
RS485, RS232 e I2C, bem como a necessidade de gerir comunicações
concorrentes, diferenças de taxa e prioridades. A capacidade nominal
mínima da interface Ethernet é de 100~Mbit/s.

Distingue-se este núcleo obrigatório das características desejáveis e
das decisões ainda abertas. Os protocolos, as configurações, a escala
de utilização e os critérios de aceitação pendentes serão concretizados
através das questões da Secção~\ref{sec:questoes-validacao}. A especificação
permanece preliminar nesses aspetos, sem reabrir as capacidades já
confirmadas.

\section{Requisitos Funcionais}

Os requisitos funcionais delimitam as capacidades obrigatórias do
\textit{gateway}. RF06 e RF07 estabelecem os dois sentidos de comunicação,
e RF09 aplica-os a cada uma das quatro ligações entre Ethernet e as
restantes interfaces. A bidirecionalidade não implica transmissão
simultânea nos dois sentidos em todos os barramentos. A concorrência
entre interfaces é exigida por RF12; o número de canais e as condições
de carga permanecem por definir.

\begin{description}
    \item[RF01 --- Suporte para I2C:] \hfill\newline
    O sistema deve suportar comunicação através de uma interface I2C.

    \item[RF02 --- Suporte para RS232:] \hfill\newline
    O sistema deve suportar comunicação através de uma interface RS232.

    \item[RF03 --- Suporte para RS485:] \hfill\newline
    O sistema deve suportar comunicação através de uma interface RS485.

    \item[RF04 --- Suporte para CAN:] \hfill\newline
    O sistema deve suportar comunicação através de uma interface CAN.

    \item[RF05 --- Suporte para Ethernet:] \hfill\newline
    O sistema deve disponibilizar Ethernet como interface principal de
    comunicação com o computador ou sistema central, com capacidade
    nominal mínima de 100~Mbit/s. Este valor não constitui uma garantia
    de débito útil ou de latência.

    \item[RF06 --- Escrita e transmissão:] \hfill\newline
    O sistema deve permitir operações de escrita ou transmissão de dados
    através de cada uma das cinco interfaces previstas, de acordo com as
    características da respetiva tecnologia.

    \item[RF07 --- Leitura e receção:] \hfill\newline
    O sistema deve permitir operações de leitura ou receção de dados
    através de cada uma das cinco interfaces previstas, de acordo com as
    características da respetiva tecnologia.

    \item[RF08 --- Processamento e adaptação:] \hfill\newline
    O sistema deve processar a informação recebida e efetuar as adaptações
    necessárias à operação de destino, de acordo com as configurações e
    os protocolos envolvidos. As operações suportadas devem explicitar
    os dados de entrada, as transformações e os resultados esperados.

    \item[RF09 --- Encaminhamento bidirecional através de Ethernet:] \hfill\newline
    O sistema deve permitir ao sistema central enviar comandos ou dados
    e receber dados ou feedback através de Ethernet, de e para dispositivos
    ligados a cada uma das interfaces CAN, RS485, RS232 e I2C. Cada operação
    requer informação suficiente sobre origem, destino, configuração e
    protocolo, incluindo o mapeamento ou a transformação quando necessários.
    A operação só pode ser executada se essa informação permitir determinar
    uma ação compatível com a interface e o protocolo de destino.

    \item[RF10 --- Preservação da informação necessária:] \hfill\newline
    O processamento e o encaminhamento devem preservar os dados e metadados
    necessários à operação da interface ou protocolo de destino, mantendo
    o seu significado de acordo com o mapeamento definido. Esses elementos
    e as transformações permitidas devem ser identificados para cada
    operação, permitindo comparar a informação de origem com o resultado
    no destino.

    \item[RF11 --- Acesso externo e seleção do destino:] \hfill\newline
    Através de Ethernet, o sistema deve permitir solicitar uma operação,
    indicando a interface ou barramento de destino, o dispositivo, endereço
    ou identificador quando aplicável, e os dados ou comandos a transmitir.
    Os dados e feedback devolvidos ao sistema central devem permitir
    identificar a respetiva origem.

    \item[RF12 --- Comunicações concorrentes:] \hfill\newline
    O sistema deve manter diferentes interfaces ativas e gerir operações
    de comunicação concorrentes, incluindo receção em várias interfaces.
    A espera por uma interface não deve impedir o progresso das operações
    nas restantes, nas condições de funcionamento e carga a definir.

    \item[RF13 --- Gestão de diferenças de taxa:] \hfill\newline
    O sistema deve gerir diferenças entre as taxas de entrada e saída,
    assegurando armazenamento temporário da informação que aguarda
    transmissão quando necessário. Os limites de carga, duração da
    sobrecarga e capacidade de armazenamento, bem como o comportamento
    quando forem excedidos, serão definidos através de QV10.

    \item[RF14 --- Gestão de prioridades:] \hfill\newline
    Quando existam comunicações concorrentes a disputar recursos, o sistema
    deve gerir o seu atendimento segundo uma política de prioridades a
    definir. A política concreta e a possibilidade de a configurar durante
    a utilização permanecem por validar.
\end{description}

O encaminhamento obrigatório abrange todos os quatro pares Ethernet
com CAN, RS485, RS232 e I2C; não se reduz à escolha de alguns desses pares
para demonstração. Os casos concretos e os protocolos suportados ainda
terão de ser definidos. Não se exige conversão semântica automática
entre protocolos desconhecidos, nem correspondência de uma trama de
entrada para uma trama de saída. Em CAN, um identificador de trama não
constitui, por si só, um endereço de dispositivo; a seleção deve respeitar
o significado atribuído pelo protocolo utilizado.

A generalização a pares sem Ethernet é uma extensão desejável, não um
critério obrigatório de aceitação. RF11 especifica o acesso externo;
a uniformidade e a transparência desse acesso são objetivos desejáveis
de abstração, sem impor uma representação intermédia universal ou um
formato interno.
O armazenamento temporário exigido em RF13 não garante absorção ilimitada
de tráfego superior à capacidade de saída.

\section{Requisitos Não Funcionais}

As propostas RNF01, RNF02 e RNF03 dizem respeito à organização e evolução
do sistema. Os restantes requisitos descrevem propriedades de fiabilidade,
observabilidade e comportamento temporal decorrentes do âmbito funcional. Os critérios de
aceitação ainda não confirmados serão validados com os docentes; os
métodos previstos encontram-se na secção de rastreabilidade.

\begin{description}
    \item[RNF01 --- Modularidade:] \hfill\newline
    A organização do sistema deverá distinguir as responsabilidades
    associadas a cada interface das funções de processamento e
    encaminhamento, com fronteiras e dependências identificadas
    na arquitetura e na implementação.

    \item[RNF02 --- Manutenção:] \hfill\newline
    A documentação deverá identificar as responsabilidades, as entradas,
    as saídas e as dependências das funções implementadas, bem como as
    configurações necessárias para reproduzir as operações de comunicação.

    \item[RNF03 --- Extensibilidade:] \hfill\newline
    A arquitetura deverá permitir incorporar novos casos de encaminhamento
    e configuração através de pontos de extensão e dependências
    identificados, sem impor uma organização interna específica.

    \item[RNF04 --- Fiabilidade:] \hfill\newline
    Nas condições de funcionamento e carga a definir, o sistema deverá
    produzir os resultados previstos para as operações suportadas,
    incluindo quando existam comunicações concorrentes e diferenças
    de taxa entre interfaces.

    \item[RNF05 --- Observabilidade e diagnóstico:] \hfill\newline
    O sistema deverá permitir observar as mensagens tratadas, os resultados
    das operações e os erros detetados, associando-os às interfaces e às
    operações correspondentes. A forma de acesso a essa informação
    permanece por definir; não se pressupõe uma interface de diagnóstico
    específica nem registo persistente de todas as mensagens.

    \item[RNF06 --- Comportamento temporal:] \hfill\newline
    O desempenho temporal do gateway deverá ser caracterizado, com o
    objetivo de minimizar a latência introduzida pelo processamento e
    encaminhamento. Este requisito é preliminar e requer quantificação:
    os limites máximos aceitáveis sob carga individual e concorrente serão
    definidos após a resolução de QV07. Não existe atualmente uma garantia
    de tempo real estrito (\textit{hard real-time}).
\end{description}

\section{Rastreabilidade dos Requisitos}

As Tabelas~\ref{tab:rastreabilidade-requisitos}
e~\ref{tab:rastreabilidade-rnf} relacionam as capacidades e propriedades
com os objetivos, cenários, conclusões do Estado da Arte e validação de
requisitos com os docentes. As formas de validação são previstas: antes dos ensaios,
será necessário definir configurações, dados e resultados esperados.

\begin{table}[htbp]
    \centering
    \small
    \setlength{\tabcolsep}{4pt}
    \renewcommand{\arraystretch}{1.12}
    \caption{Origem e validação prevista dos requisitos funcionais.}
    \label{tab:rastreabilidade-requisitos}
    \begin{tabularx}{\textwidth}{@{}>{\raggedright\arraybackslash}p{0.16\textwidth}*{2}{>{\raggedright\arraybackslash}X}@{}}
        \toprule
        Requisito & Origem / Justificação & Forma prevista de validação \\
        \midrule
        RF01--RF04
        & Objetivos~\ref{obj:estudo-tecnologias} e~\ref{obj:comunicacao};
        CA01: acesso aos quatro barramentos.
        & Teste de comunicação por interface nas configurações acordadas. \\
        \addlinespace
        RF05
        & Validação com os docentes: Ethernet principal com mínimo nominal de 100~Mbit/s;
        CA01.
        & Inspeção da capacidade documentada e verificação da ligação
        a pelo menos 100~Mbit/s com equipamento compatível. \\
        \addlinespace
        RF06--RF07
        & Objetivo~\ref{obj:comunicacao}; CA01; validação com os docentes:
        comandos e feedback nos dois sentidos.
        & Testes de escrita/transmissão e leitura/receção nas cinco
        interfaces, sem exigir simultaneidade dos dois sentidos. \\
        \addlinespace
        RF08, RF10
        & Objetivo~\ref{obj:processamento-encaminhamento}; síntese do
        Estado da Arte: adaptação e preservação de significado.
        & Comparação de entradas e saídas, incluindo dados e metadados
        necessários, para as adaptações especificadas. \\
        \addlinespace
        RF09
        & Validação com os docentes; CA01: Ethernet com cada uma das outras quatro interfaces.
        & Testes nos dois sentidos de cada um dos quatro pares;
        casos com informação válida e insuficiente para executar a operação. \\
        \addlinespace
        RF11
        & Validação com os docentes; CA01: acesso externo e seleção do destino.
        & Solicitação de operações por Ethernet variando interface e
        destino; verificação da origem do feedback. \\
        \addlinespace
        RF12, RF14
        & Validação com os docentes; CA03; Estado da Arte: concorrência e prioridades.
        & Tráfego concorrente; verificação do progresso das operações e
        da ordem de atendimento segundo a política acordada. \\
        \addlinespace
        RF13
        & Validação com os docentes; Estado da Arte: diferenças de taxa e armazenamento
        temporário.
        & Entrada temporariamente mais rápida que a saída; verificação
        do armazenamento e do comportamento nos limites a definir em QV10. \\
        \bottomrule
    \end{tabularx}
\end{table}

\begin{table}[htbp]
    \centering
    \small
    \setlength{\tabcolsep}{4pt}
    \renewcommand{\arraystretch}{1.12}
    \caption{Origem e validação prevista dos requisitos não funcionais.}
    \label{tab:rastreabilidade-rnf}
    \begin{tabularx}{\textwidth}{@{}>{\raggedright\arraybackslash}p{0.16\textwidth}*{2}{>{\raggedright\arraybackslash}X}@{}}
        \toprule
        Requisito & Origem / Justificação & Forma prevista de validação \\
        \midrule
        RNF01, RNF03
        & Objetivo~\ref{obj:arquitetura}; heterogeneidade das interfaces e
        generalização desejável do encaminhamento.
        & Inspeção da arquitetura e do software; análise de um caso de
        extensão acordado e das dependências alteradas. \\
        \addlinespace
        RNF02
        & Objetivo~\ref{obj:arquitetura}; problema de manutenção
        identificado no enquadramento.
        & Inspeção da documentação e da correspondência com funções e
        configurações implementadas. \\
        \addlinespace
        RNF04
        & Objetivo~\ref{obj:validacao}; CA03; RF12--RF13:
        funcionamento sob concorrência e diferenças de taxa.
        & Ensaios repetidos; comparação de resultados esperados e
        observados. Cargas, falhas, repetições e aceitação a validar. \\
        \addlinespace
        RNF05
        & Objetivo~\ref{obj:validacao}; CA03:
        observação de mensagens, operações e erros.
        & Observação de operações e de erros provocados em condições
        controladas; verificação da sua associação às interfaces. \\
        \addlinespace
        RNF06
        & Validação com os docentes: redução de atrasos; CA03 e RF12--RF14.
        & Medição de latência com carga individual e concorrente;
        pontos de medição e limites de aceitação a definir. \\
        \bottomrule
    \end{tabularx}
\end{table}

\section{Características Desejáveis}

São desejáveis o suporte Gigabit Ethernet, a alimentação por Ethernet
(\textit{Power over Ethernet}, PoE) para instalações com um único cabo e
um consumo reduzido. Não constituem critérios obrigatórios de aceitação
nesta fase; não existe um limite de consumo confirmado. A generalização
do encaminhamento a pares sem Ethernet enquadra-se igualmente como
extensão. Estas características serão avaliadas sem presumir a sua
inclusão na primeira versão.

\section{Restrições do Projeto}

O âmbito funcional confirmado inclui o acesso bidirecional por Ethernet
aos quatro tipos de interfaces, a gestão de concorrência e a capacidade
nominal Ethernet de pelo menos 100~Mbit/s. Não estão confirmadas
restrições de implementação relativas à plataforma, componentes,
orçamento, PCB ou condições especiais de instalação.

Os protocolos, as quantidades de canais e nós, as políticas de atendimento
e os limites de armazenamento e de desempenho continuam por definir.
A necessidade de isolamento e proteção deve ser avaliada por interface;
não se impõe isolamento galvânico universal. Também não se exigem SPI,
um sistema operativo específico ou funcionamento em pressão, óleo ou
outras condições ambientais especiais. Decisões em aberto não são
restrições técnicas já estabelecidas.

\section{Requisitos a Validar com os Docentes}
\label{sec:questoes-validacao}

Estão definidos os quatro pares base através de Ethernet, os dois sentidos
de comunicação, a necessidade de concorrência e o mínimo nominal de
100~Mbit/s. As questões seguintes completam a especificação antes de a
congelar; não colocam novamente essas capacidades em dúvida.

\begin{description}
    \item[QV01 --- Operações de comunicação:] \hfill\newline
    Para cada par Ethernet com CAN, RS485, RS232 e I2C, que operações,
    papéis, configurações e adaptações devem ser demonstrados? Que
    generalizações a pares sem Ethernet se pretende estudar como extensão?

    \item[QV02 --- Protocolos de aplicação:] \hfill\newline
    Que protocolos de aplicação devem ser suportados em RS232 e RS485?
    Existem protocolos superiores obrigatórios nas restantes interfaces,
    por exemplo em CAN? Como serão interpretados endereços e identificadores?

    \item[QV03 --- Plataforma de execução:] \hfill\newline
    Existem restrições à utilização de MCU, SoC ou plataforma com Linux?
    Os requisitos temporais a definir exigem um RTOS ou deixam em aberto
    o mecanismo de execução?

    \item[QV04 --- Variante de CAN:] \hfill\newline
    O suporte para CAN deve abranger Classical CAN, CAN FD ou ambas
    as variantes?

    \item[QV05 --- Acesso por Ethernet:] \hfill\newline
    Que protocolos de transporte e de aplicação serão utilizados: TCP,
    UDP ou outros? Qual será o formato externo para solicitar operações,
    selecionar interface e destino e receber dados, feedback e erros?

    \item[QV06 --- Placa de circuito impresso:] \hfill\newline
    É necessário desenvolver uma placa de circuito impresso
    (PCB) própria?

    \item[QV07 --- Desempenho e fiabilidade:] \hfill\newline
    Além do mínimo nominal Ethernet de 100~Mbit/s já definido, quais são
    os limites aceitáveis de latência e as cargas e condições de operação?
    Que falhas, repetições e critérios serão usados na avaliação de
    fiabilidade? Não existe ainda uma exigência de tempo real estrito.

    \item[QV08 --- Quantidades de interfaces e dispositivos:] \hfill\newline
    Quantos canais físicos de cada tipo devem existir, incluindo eventuais
    interfaces repetidas? Quantos nós ou dispositivos devem ser suportados
    por barramento, em particular CAN e RS485? Que combinações devem operar
    simultaneamente?

    \item[QV09 --- Prioridades:] \hfill\newline
    Qual será a política de atendimento das comunicações concorrentes?
    Que critérios de prioridade e garantias de progresso devem ser
    observados? As prioridades terão de ser configuráveis durante a utilização?

    \item[QV10 --- Armazenamento temporário:] \hfill\newline
    Que diferenças de taxa e durações de tráfego devem ser absorvidas?
    Quais serão os limites de armazenamento e o comportamento quando
    forem excedidos? Que mecanismos de controlo de fluxo permitem os
    protocolos selecionados?

    \item[QV11 --- Isolamento e proteção:] \hfill\newline
    Que interfaces necessitam de isolamento galvânico e que condições
    elétricas, exposição a ruído, surtos ou descargas eletrostáticas
    devem orientar a proteção? Quais serão os critérios de avaliação?

    \item[QV12 --- Características opcionais:] \hfill\newline
    Que características desejáveis se pretende incluir na primeira versão:
    Gigabit Ethernet, PoE ou ambas? Que condições
    e eventual limite de consumo devem orientar a avaliação energética?
\end{description}
